\documentclass[10pt,a4paper]{amsart}
\usepackage[margin=19mm]{geometry}
\usepackage{amsmath,amssymb,mathtools}
\usepackage[scr=rsfs,cal=euler]{mathalfa}
\usepackage{xfrac}
\usepackage[unicode,hidelinks,bookmarks=false]{hyperref}
\newcommand{\ie}{i.e.\ }
\newcommand{\cf}{cf.\ }
\newcommand{\resp}{respectively\ }
\newcommand{\Subsection}{\subsection}
\providecommand{\nobreakdash}{\nobreak\hbox{-}}
\setlength{\emergencystretch}{2em}
\multlinegap0pt
\usepackage{mathrsfs}
\usepackage{bbm}
\usepackage{url}
\usepackage{bbm}
\usepackage{eucal}
\usepackage{tikz-cd}
\usepackage{stmaryrd}
\newtheorem{Theo}{Theorem}[section]
\newtheorem{Def}[Theo]{Definition}
\newtheorem{Ex}[Theo]{Example}
\newtheorem{Example}[Theo]{Example}
\newtheorem{Prop}[Theo]{Proposition}
\newtheorem{Lemma}[Theo]{Lemma}
\newtheorem{Corolary}[Theo]{Corollary}
\newtheorem{Rem}[Theo]{Remark}
\newtheorem*{CH}{Continuum Hypothesis (CH)}
\newcommand{\C}{\mathcal{C}}
\newcommand{\Sets}{\mathcal{S}\text{ets}}
\newcommand{\E}{\mathcal{E}}
\newcommand{\One}{\mathbbm{1}}
\newcommand{\Po}{\mathbb{P}}
\newcommand{\B}{\mathbb{B}}
\newcommand{\He}{\mathbb{H}}
\newcommand{\Qo}{\mathbb{Q}}
\newcommand{\D}{\mathcal{D}}
\newcommand{\I}{\mathrm{I}}
\newcommand{\N}{\omega}
\newcommand{\Z}{\mathbb{Z}}
\newcommand{\Pk}{\mathbb{P}_{\kappa}}
\newcommand{\Tops}{\mathcal{T}ops}
\newcommand{\hC}{\Hat{\C}}
\newcommand{\Dnk}{D_{(n,k)}}
\newcommand{\Ekl}{E_{(k,l)}}
\newcommand{\ydot}{\overset{.}{y}}
\newcommand{\zdot}{\overset{.}{z}}
\newcommand{\xdot}{\overset{.}{x}}
\newcommand{\wdot}{\overset{.}{w}}
\newcommand{\forces}{\Vdash^{*}}
\newcommand{\Forces}{\Vdash}
\newcommand{\lb}{\llbracket}
\newcommand{\rb}{\rrbracket}
\newcommand{\us}{\overset{\to}{u}}
\newcommand{\lp}{\overset{\leftarrow}{p}}
\newcommand{\x}{\overset{\to}{x}}
\usepackage{setspace}
\begin{document}
\title[Relating forcing relations]{Relating forcing relations}
\author{Michel Viana Smykalla; TALTECH; Hugo Luiz Mariano; IME-USP}
\date{}
\maketitle
\noindent\emph{Source and licence.} Relating forcing relations. Version arXiv:2605.26896v1 (CC BY 4.0). This material is distributed under the Creative Commons Attribution 4.0 International licence, http://creativecommons.org/licenses/by/4.0/, as recorded in the arXiv OAI licence metadata for this exact version and shown on the abstract page https://arxiv.org/abs/2605.26896v1. Full rights notice: licenses/arxiv-2605.26896-CC-BY-4.0.txt.\par\smallskip
\noindent\textit{Editorial scope.} Counted unit: the original \texttt{Lemma} environment \texttt{BeilinsonLemma} at source lines 960--981 together with its complete proof at lines 983--985; the delivered document keeps source lines 779--1050 so that every referenced label and every citation resolves inside it. Native editable source is the paper's own concatenated TeX; the AMS wrapper is editorial formatting only. No publisher class, upstream script or unrelated artwork is redistributed.
\section{Forcing, Boolean-valued models and sheaves over Boolean algebras}\label{forcingmodelsandsheaves}



Let $\B$ be a complete Boolean algebra. It is well known (see, for instance, \cite{Bell05} Appendix) that from the Boolean-valued model $V^\B$ can be extracted a category  $Set^{(\B)}$, "by taking quotients". In more detail:
\begin{itemize}
\item An object of  $Set^{(\B)}$ is a class of equivalence $[x]$, where $x \in V^\B$ and $[x] = [x']$ iff $||x = x'|| = 1_{\B}$;
\item An arrow $[f] : [x] \to [y]$ in $Set^{(\B)}$ is a class of equivalences $[f]$, where $f \in V^\B$,  $[f] = [f']$ iff $||f = f'|| = 1_{\B}$, and $|| f$ is a function with domain $x$ and range contained in $y|| = 1_\B$. 
\end{itemize}

Moreover, in Appendix  of \cite{Bell05}, it is sketched an equivalence of categories between this category obtained by quotients on the Boolean-valued universe $V^\B$ and the category of sheaves of sets on $\B$ obtained from the natural notion of covering given by suprema: $$Set^{(\B)} \simeq Sh_{\bigvee(\B)}$$.



    



 


 
% \section{Equivalence of categories induced by a dense morphism}\label{equivalenceofcategoriesdensemorphism}
%Escreve na velocidade da luz kct
%Ao infinito e alem!!!

Taking into account the above described  scenario, it is possible to establish a connection between the sheaf theoretic version of forcing and Boolean-valued models by showing an equivalence of categories. In this section, we will discuss this result as well as a generalization of it.

We will assume basic knowledge in category theory, introducing the concepts of Grothendieck topology and Grothendieck topos. Given a category $\C$, $Obj(\C)$ represents the (class of) objects of $\C$. As a reference for this introduction to Topos Theory, see \cite{MacMoe94}.


\begin{Def}\label{Sievedef}
	Let $\C$ be a category and $C \in \text{Obj}(\C)$. A \textbf{sieve} on $C$ is a family $S$
	of arrows in $\C$ all with codomain $C$, such that if $f: A \to C$ belongs to $S$ and $g: B \to A$ is any arrow in $\C$ with codomain $A$, then $f \circ g \in S$.
	
\end{Def}

If $\C$ is a locally small category, a sieve on an object $C$ of $\C$ will be a subobject $S$ of $y(C) = Hom_{\C}( - ,C)$, i.e., there exists a monomorphism $S \rightarrowtail y(C)$ in $\hC$ satisfying the universal property of subobjects of $y(C)$. Moreover, if $S$ is a sieve on $C$ and $h: D \to C$ is any arrow in $\C$ with codomain $C$, then 
$$h^{*}(S) = \{g\text{ }|\text{  codom}(g) = D \text{ and } h \circ g \in S\}$$
is a sieve on $D$.

\begin{Def}
	Let $\C$ be a small category. A \textbf{Grothendieck topology} on $\C$ is a function $J$ which associates to each object $C$ of $\C$ a family $J(C)$ of sieves on $C$ satisfying the following properties.
	\begin{enumerate}
		\item $t_{C} = \{ f\text{ } | \text{  codom}(f) = C\} \in J(C)$. We call $t_{C}$ the \textbf{maximal sieve}.
		
		\item If $S \in J(C)$, then for any arrow $h: D \to C$ in $\C$, $h^{*}(S) \in J(D).$ This property is known as the \textbf{stability axiom}.
		
		\item If $S \in J(C)$ and $R$ is a sieve on $C$ such that for any $h: D \to C \in S$, $h^{*}(R) \in J(D)$, then $R \in J(C)$. Some books call this condition by the \textbf{transitivity axiom}. 
	\end{enumerate}
\end{Def}

\begin{Def}\label{Sitedef}
We call by \textbf{site} a pair $(\C,J)$, composed by a small category $\C$ and a Grothendieck topology $J$ on $\C$. If $S \in J(C)$, we say that $S$ \textbf{covers} $C$. 
\end{Def}

Using a Grothendieck topology we will construct the category where the $CH$ is not satisfied. The next example is the topology that in fact will provide the Cohen topos.

\begin{Example}
	Let $\Po$ be a poset. Note that $\Po$ is a category, where an arrow $p \to q$ in $~\Po$ means that $ p\leq q$. Given $p \in \Po$, consider the set $\overset{\leftarrow}{p} = \{q \in \Po\text{ } | \text{ }q \leq p\}$. A subset $D \subseteq \overset{\leftarrow}{p}$ is \textbf{dense below $p$} if for every $r \leq p$, there exists $q \in D$ such that $q \leq r$. The dense sieves form a Grothendieck topology $J$ on $\Po$ by
	
	$$J(p) = \{D : D \text{ is a sieve on } p \text{ and dense below p}\}. $$
	We call this Grothendieck topology by \textbf{dense topology} or \textbf{double-negation topology}, denoted by $\neg\neg$-topology or just $\neg\neg$.  
\end{Example}
 A generalization of Boolean algebras is Heyting algebras, where we remove the \emph{excluded middle}, that is, $-(-u) = u$ is not true in  a Heyting algebra.
\begin{Example}
	Let $H$ be a complete Heyting algebra. We can see $H$  as a category in the same way as posets (an arrow $h \to k$ in $H$ means that $h \leq k$, for $h,k \in H$). The \textbf{sup topology} on $H$ is a Grothendieck topology $J$ such that for all $h \in H$,
	
	$$J(h) = \{S: \bigvee  S = h\}.$$
We usually denote the sup topology by $\bigvee$.
\end{Example}

\begin{Def}
	Let $\C$ be a category with pullbacks. A \textbf{basis} for a Grothendieck topology on $\C$ is a function $B$ which associates an object $C \in \text{Obj}(\C)$ to a collection $B(C)$ of families of arrows in $\C$ with codomain $C$ satisfying the next conditions:
	
	\begin{enumerate}
		\item If $f: C' \to C$ is an iso, then $~\{f:C' \to C\} \in B(C)$.
		
		\item If $~\{f_{i}: C_{i} \to C ~ |\text{ for all }i \in I\} \in B(C)$ and $g: D
		 \to C$ is any morphism of $\C$ with codomain $C$, then the family of projections $\{ \pi_{2}: C_{i} \times D \to D~|\text{ for all } i \in I\}$ belongs to $B(D)$.
		 
		\item If $~\{f_{i}: C_{i} \to C ~ |\text{ for all }i \in I\} \in B(C)$ and for each $i \in I$, there exists a family $\{g_{ij}: D_{ij} \to C_{i}~|~ j \in J_{i}\} \in B(C_{i})$, then $\{f_{i} \circ g_{ij}: D_{ij} \to C ~|\text{ for all }i \in I, j \in J_{i}\} \in B(C)$.
	\end{enumerate}
\end{Def}
\begin{Def}
	Let $\C$ be a category and $A \in \text{Obj}(\C)$. Given a collection of morphism $B = \{g_{i}:A_{i} \to A\}_{ i \in I }$ in $\C$, the sieve (see Definition \ref{Sievedef}) \textbf{generated} by $B$, is the sieve $$<B> = \{f \circ g_{j}: f: C \to A_{j} \text{ and } g_{j} \in B\}.$$ 
\end{Def}

Given a basis $B$ on $\C$, we can obtain a Grothendieck topology $J$ defining for each $C \in \text{Obj}(C)$, for every sieve $S$ on $C$,
$$S \in J(C) \text{ if and only if there exists } R \in B(C)\text{ such that }<R> \subseteq S. $$


The objects of the category that we are looking for (Grothendieck topos) are sheaves. To talk about them, we need more definitions, starting with the notion of presheaves. 

\begin{Def}\label{presheafdef}
	Let $\C$ be a category. A \textbf{presheaf} $F$ on $\C$ is a functor
	$$F: \C^{op} \to \Sets.$$
	
\end{Def}

\begin{Example}\label{Representablepresheaf}
One example of a family of presheaves are the functors of the form $y(C)=Hom_{\C}( - ,C)$ for a given object $C$ of $\C$. We call them by \textbf{representable presheaves}.
\end{Example}

\begin{Def}\label{categoryofpresheavesdef}
	
	Let $\C$ be a category. The \textbf{category of presheaves on} $\C$ is the category ${\Sets}^{\C^{op}}$ of functors $F: \C^{op} \to \Sets$ from the opposite category of $\C$ to $\Sets$ and natural transformations between them. In this case, we represent this category of presheaves over $\C$ using the notation $\Hat{\C}$. 
	
\end{Def}

Now, fix a site $(\C,J)$ and a presheaf $P$ on $\C$. Given an arrow $g: E \to D$ in $\C$, we have that $P(g) : P(D) \to P(E)$. For all $x \in P(D)$, $x \cdot g$ stands for $P(g)(x)$.  If $S$ is sieve and covers an object $C$ of $\C$, a \textbf{matching family} for $S$ of elements of $P$ is a function which associates each element $f: D \to C$ of $S$, to an element $x_{f} \in P(D)$ satisfying $$P(g)(x_{f}) = x_{f} \cdot g = x_{f\circ g},$$ for all morphisms $g: E \to D$ of $\C$. An \textbf{amalgamation} of such a matching family is a single element $x \in P(C)$ such that $$x \cdot f = x_{f} \text{ for all } f \in S.$$

\begin{Def}\label{Sheafdef}
	Let $(\C,J)$ be a site and $P$ a presheaf over $\C$. Then $P$ is a \textbf{sheaf for} $J$ if for every matching family  of elements of $P$ for any cover of any object of $\C$ there exists a unique amalgamation. In this case, we also say that $P$ is a sheaf on the site $(\C,J$).
\end{Def}

Then, sheaves on a site $(\C,J)$ form a category $Sh(\C,J)$, where the objects are the sheaves and the arrows, natural transformations between them. In this case, $Sh(\C, J)$ is a full subcategory of $\Sets^{\C^{op}}$, then we have the inclusion functor

%\shorthandoff{"}
\begin{center}
	% https://tikzcd.yichuanshen.de/#N4Igdg9gJgpgziAXAbVABwnAlgFyxMJZABgBpiBdUkANwEMAbAVxiRAGUALACgB1eAwqQBSAShABfUuky58hFAEZyVWoxZt+7GDjgA9YPwEGIaCRMmqYUAObwioAGYAnCAFskZEDghJl3uiwGNk4ICABrEGp6ZlZEECxLCSA
	\begin{tikzcd}
		{Sh(\C,J)} \arrow[r, "Id_{Sh(\C,J)}", hook] & \Sets^{\C^{op}}.
	\end{tikzcd}
\end{center}

\begin{Def}
	A \textbf{Grothendieck topos} is a category which is equivalent to the category $Sh(\C,J)$ of sheaves on some site $(\C,J)$.
\end{Def}


\begin{Ex}
	The following categories are Grothendieck toposes and we will use them later:
	\begin{enumerate}
		\item The category of sets, $\Sets$.
		\item The category of presheaves $~\widehat{\C} = \Sets^{\C^{op}}$, where $\C$ is a small category.
	\end{enumerate}
\end{Ex}
%----------------


\begin{Def}\label{separativeposet}
    A poset $\Po$ is called separative if for all $p, q \in \Po$, if $p \nleq q$, then there exists $r \leq p$ such that for all $s \leq r$, holds $s \nleq q$.
\end{Def}


 Many relevant posets that appear in  forcing are separative. In particular the so called Cohen forcing is separative and, for any complete Boolean algebra $\B$, the poset $\Po = \B \setminus \{0\}$ is separative (just take $r = p . (- q)$).

The (well known) result concerning separative forcing is the following:



 
    
    


\begin{Theo}\label{equivalenceofcategoriesdensemorphismseparative}
	Let $\Po$ be a separative poset. Then there exists a complete Boolean algebra $\mathbb{B}$ such that 
	$$Sh_{\neg \neg}(\Po) \simeq Sh_{\bigvee}(\mathbb{B}).$$
\end{Theo}
\begin{proof}
See \cite{MacMoe94},  Corollary 3 of the Section 4 of the Appendix: Sites for Topoi, and the following comments, pages 590-591.
\end{proof}

One would interpret this theorem as the separative forcing, as a method, has the \emph{same} content as Boolean-valued models. 

Now we will give a generalization of Theorem \ref{equivalenceofcategoriesdensemorphismseparative} for any forcing poset as part of the original contributions of this work (Theorem \ref{HugoMicheltheorem2}).




\begin{Def}
	Let $\C$ be a small category . Let $J$ and $K$ be two Grothendieck topologies on $\C$. We say that $J$ is \textbf{finer} than $K$ if for all $C \in \text{Obj}(\C)$, $J(C)\subseteq K(C)$.
\end{Def}

\begin{Def}
	Let $\C$ and $\D$ be a small categories and $(\D,J)$ a site. Given a functor $F: \C \to \D$, the \textbf{topology induced} by $F$ on $\C$, denoted by $J_{F}$, is the finest one such that for all $H$ sheaf on $\D$ for $J$, $H \circ F$ is a sheaf on $\C$ for $J_{F}$.
\end{Def}



\begin{Lemma}\label{BeilinsonLemma}
	Let $\C$ and $\D$ be  small categories and $J$ a Grothendieck topology on $\D$. Suppose that there exists a faithful functor $F: \C \to \D$ satisfying the following:
	\begin{enumerate}
		
		
		\item[] Fix an object $D \in Obj(\D)$. Given a finite set of objects $\{C_{i} \in Obj(\C): i \in I\}$ and a family of morphism in $\D$ of the form $(D \to F(C_{i}))_{i \in I}$, there exists a family of morphisms  in $\D$ of the form $(F(E_{j}) \to D)_{j \in J}$ such that the composition $F(E_{j}) \to D \to F(C_{i})$ lies in the image of $Hom(E_{j},C_{i}) \to Hom_{\D}(F(E_{j}),F(C_{i}))$, for all $i \in I, j \in J$, and the sieve generated by $(F(E_{j}) \to D)_{j \in J}$ covers $D$.
	\end{enumerate}
	Then:
	\begin{enumerate}
		\item The induced topology $J_{F}$ on $\C$ has the following property: For all $C \in Obj(\C)$,
		$$S \in J_{F}(C) \text{ iff } <F(S)> \in J(F(C)),$$
		where $<F(S)>$ is the sieve generated by $ \{F(f): f \in S\}$.
		
		\item $Sh_{J_{F}}(\D) \simeq Sh_{J}(\C)$, i.e., the functor
		\begin{align*}
				F^{*}:\widehat{\D} &\to \widehat{\C} \\
							H &\mapsto H \circ F
		\end{align*}
		restricts to the equivalence of categories above.
	\end{enumerate}
	
\end{Lemma}

\begin{proof}
	See \cite{Bei12} first Proposition of Section 2.1.
\end{proof}


A corollary of Lemma \ref{BeilinsonLemma} is the well-known Comparison lemma:

\begin{Corolary}\label{comparisonlemma}
	Let $\C$ be a small category and $(\D,J)$ a site. Let $F: \C \to D$ be a full and faithful functor. Let $J_{F}$ be the induced topology on $\C$ by $F$. If for all $d \in \text{Obj}(\D)$ there exists a sieve generated by morphisms of the form $(F(C_{i}) \to d)_{i \in I})$, for $C_{i} \in \text{Obj}(\C)$, then
	$$Sh_{J_{F}}(\C) \simeq Sh_{J}(\D).$$
\end{Corolary}

\begin{proof}
	See Lemma \ref{BeilinsonLemma}.
\end{proof}

\begin{Prop}
    \label{sheavesposetalgebrasem0}
	Let $\Po$ be a forcing poset and $\B$ its complete Boolean algebra of regular open sets of $~\Po$. Then
	$$Sh_{\neg \neg}(\Po) \simeq Sh_{\neg \neg}(\B \setminus \{ \emptyset\}).$$
\end{Prop}

\begin{proof}
	Let $i: \Po \to \B$ be the canonical dense morphism. To use Lemma \ref{BeilinsonLemma}, we need to check that the dense morphism $i: \Po \to \B \setminus \{\emptyset\}$ satisfies the required property. So fix $b \in \B \setminus \{\emptyset\}$. Let $$(b \leq  i(a_{j}))_{j \in J}$$ be a finite family of morphism in $\B \setminus \{\emptyset\}$. By density, there exists $a_{b} \in \Po$ such that $i(a_{b})  \leq b$. From $i(a_{b}) \leq b \leq i(a_{j})$, we conclude that $i(a_{b}) \not \perp i(a_{j})$, for all $j \in J$.Then, for each $j \in J$, there exists $r_{j} \in \Po$ such that $r_{j} \leq a_{b}, a_{j}$. Let $S$ be the sieve generated by $(i(r_{j}) \leq b)_{j \in J}$. We need to show that $S \in \neg \neg (b)$. Let $q \leq b$. Fix $j' \in J$. By density, there exists $r_{q} \in \Po$ such that $i(r_{q}) \leq q \cdot i(r_{j'})$. Then $i(r_{q}) \leq q$ and $i(r_{q}) \in S$, because $i(r_{q}) \leq i(r_{j'})$. Note that for all $j \in J$, the composition $i(r_{j}) \leq b \leq i(a_{j})$ lies on the image of 
	$$Hom_{\Po}(r_{j},a_{j}) \to Hom_{\B \setminus \{\emptyset\}}.(i(r_{j}),i(a_{j})).$$
	
	Consider $\neg \neg$ the double-negation topology on $\B \setminus \{\emptyset\}$. By Lemma \ref{BeilinsonLemma}, the induced topology $J_{i}$ provides an equivalence of categories of sheaves
	
	$$Sh_{J_{i}}(\Po) \simeq Sh_{\neg \neg}(\B \setminus \{ \emptyset \}).$$
	
	We now show that $J_{e}$ is the double negation topology $\neg \neg$. Fix a $p \in \Po$ and let $$S = (p_{j} \leq p)_{j \in J} \in \neg \neg(p).$$ We will show that  $<i(S)> \in \neg \neg (i(p))$. Let $b \in \B \setminus \{\emptyset\}$ such that $b \leq i(p)$. By density, there exists $a \in \Po$ such that $i(a) \leq b$. In particular, $i(a)\not \perp i(p)$, so there exists $c \in \Po$ such that $c \leq a,p	$. But $S$ is a sieve, then $ (c \leq p) \in S$. We conclude that $(i(c)\leq i(p)) \in <i(S)>$ and $i(c) \leq b.$
	
	Conversely, suppose that $<i(S)> \in \neg \neg (i(p))$. Fix a $q \leq p$. In particular, $i(q) \leq i(p)$. Then, there exists $(s \leq i(p)) \in <i(S)>$ so that $s \leq i(q)$. Moreover, there exists $u \in S$ such that $i(u) \leq s$. From $i(u) \not \perp i(q)$, we conclude that there is $v \in \Po$ such that $v \leq u,q$. Then $(v\leq i(p)) \in S$ and $v\leq q$.
\end{proof}

\begin{Prop}\label{sheavesnegationandsup}
	Let $\B$ be a complete Boolean algebra. Then
	$$Sh_{\neg \neg}(\B \setminus{\{0\}}) \simeq Sh_{\bigvee}(\B),$$
	where $\bigvee$ is the sup (Grothendieck) topology on $\B$.
\end{Prop}
\begin{proof}
	Denote by $e: \B \setminus \{ 0 \} \to \B$ the inclusion function. Note that since $e$ is the identity map, the property required to use Corollary \ref{comparisonlemma} is satisfied. Consider the sup topology $\bigvee$ on $\B$. By Corollary \ref{comparisonlemma}, the induced topology $J_{e}$ on $\B \setminus \{ 0\}$ provides an equivalence of categories of sheaves 
	$$Sh_{J_{e}}(\B \setminus \{0\}) \simeq Sh_{\bigvee}(\B).$$
	It remains to proof that $J_{e} = \neg \neg$. Fix $b \in \B \setminus \{ 0\}$ and let $S$ be a sieve that covers $b$. We will show that
	$$S \in \neg \neg(b) \text{ iff } <e(S)> \in \bigvee (e(b)).$$
	
	Suppose that $S \in \neg \neg (b)$. Denote $S$ by $ \{b_{j} \leq b\}_{j \in J}$. Then 
	$$<e(S)> = \{b_{j} \leq b\}_{j \in J} \cup \{0 \leq b\}. $$
	
	We want to show that 
	$$\bigvee \{b_{j} \leq b\}_{j \in J}\cup \{ 0 \leq b\} = \bigvee \{b_{j}\}_{j \in J} = b.$$
	
	Note that for all $j \in J$, $b_{j} \leq b$.  Therefore $\bigvee \{b_{j}\}_{j \in J} \leq b$. Suppose that it is not true that $b \leq \bigvee \{b_{j}\}_{j \in J}$. Denote by $x = \bigvee \{b_{j}\}_{j \in J}.$ Define $p = b \cdot (-x).$ Note that $p \leq b$ and $p \neq 0$. Then, there exists $j' \in J$ such that $b_{j'} \leq p$. In particular,
	$$b_{j'} \leq p \cdot (-x) \leq p \cdot (-b_{j'}).$$
	Then $b_{j'} = b_{j'}\cdot(p\cdot(-b_{j'})) = 0$. Contradiction. 
	
	On the other hand, suppose that $\bigvee\{b_{j}\}_{j \in J} = b.$ Fix $q \leq b$. Then
	$$q = q \cdot b = q \cdot(\underset{j \in J}{\bigvee}b_{j}) = \underset{j \in J}{\bigvee} q \cdot b_{j}.$$
	
	Then, there exists $j_{1} \in J$ such that $q \cdot b_{j_{2}} = b_{j_{1}}$, for some $j_{2} \in J$ (otherwise $q = 0$, which is not the case). Then $b_{j_{1}} \leq q$.
\end{proof}

\begin{Theo}\label{HugoMicheltheorem2}
	Let $\Po$ be a forcing poset. Then there exist a complete Boolean algebra $\B$ such that
	$$Sh_{\neg \neg}(\Po) \simeq Sh_{\bigvee}( \mathbb B).$$
\end{Theo}
\begin{proof}
	The result follows from Proposition \ref{sheavesposetalgebrasem0} and Proposition \ref{sheavesnegationandsup}.
\end{proof}

\begin{thebibliography}{99}
\bibitem[Bei12]{Bei12} Alexander Beilinson. p-adic periods and derived de Rham cohomology. Journal of the American Mathematical Society 25(3):715-738, 2012.
\bibitem[Bell05]{Bell05} John L. Bell. Set theory: Boolean-valued models and independence proofs. 3rd ed. Oxford University Press, 2005.
\bibitem[MacMoe94]{MacMoe94} Saunders Mac Lane and Ieke Moerdijk. Sheaves in Geometry and Logic: a first introduction to topos theory. Springer-Verlag New York, 1st edition, 1994.
\end{thebibliography}
\end{document}
